% Discussion (Gate C; deepened with literature at the user's request): the central result first; each pillar against prior work;
% the disease consequence; digital twins as experimental systems with six predictions for flies; a closing line. Novelty is
% argued against named precedents only.

A fixed connectome, trained end to end, holds working memory in neuronal adaptation and recalls it only in the brain state
that wrote it. The trace needs neither persistent firing nor a dedicated circuit, and the learned readout can use it only where
the operations that write and read it match the state of learning. The result joins two questions that are usually asked
apart: what holds a silent memory, and why recall depends on the state of the brain.

One dependence accounts for four observations in the model. The operator that turns the silent trace into a push depends on
the state of the network, in time and in excitability. Hence the lock to the trained timings, its widening when training spans
timings or excitability states, the loss of access after a lesion and its partial return on relearning. Excitability-driven
decoder drift\cellcite{haimerl2025} and state-dependent retrieval\cellcite{jovasevic2015,meyer2017} then become two sides of one
rule: recall spans the states of learning, at writing and at reading. Context-dependent memory in humans, with recall best in
the setting of learning\cellcite{smith2001}, may express the same rule at the scale of whole environments. Because the lesion
rewrites the trace as well as its reading, restoring either end alone does not restore recall.

The wiring itself is robust. Routes from the input to the readout persist after heavy cell loss, and a single extra spike is
absorbed before the decision. What fails under the lesion is therefore not reachability but the match between the state that
wrote a memory and the state that reads it. In animals, neuromodulators shift the intrinsic properties on which such a match
rests: cholinergic activation, for example, switches the firing mode of hippocampal CA3 neurons\cellcite{kis2024}. Training
across states widens the match, and cell-to-cell differences in adaptation leave the mechanism intact, in line with the
robustness that heterogeneity lends to learning in spiking networks\cellcite{pereznieves2021}. Varied conditions during learning
may therefore protect a memory against later shifts of state.

Learning on real wiring recruits adaptation, an intrinsic and cell-autonomous variable, rather than persistent firing or
synaptic facilitation\cellcite{mongillo2008,masse2019,pals2020}. Holding the cue costs no extra spikes, an economy that matters
because neural signaling is energetically expensive\cellcite{niven2008}. The time dependence of the readout also has a concrete
place. In primate recordings, stable population codes sit beside impulse responses that change with
time\cellcite{murray2017,wolff2020}; in the model, the trace persists through the delay while the operator that reads it drifts,
so the time dependence lies in the reading, not in the code. Models of representational drift ask how stable readouts can track
a code that keeps changing\cellcite{rule2022}; the model poses the converse problem, a persistent code read by a changing
operator, which the learned readout solves only over the timings it was trained on.

The carriers of the trace are enriched for \gene{SK} in the fly's own transcriptomes. SK currents shape afterhyperpolarization
and sparse coding in fly circuits\cellcite{chen2026}, and enhancing them impairs hippocampal memory in mice\cellcite{rice2024}.
The signature sides with a neurocentric view of memory\cellcite{titley2017} without implying that flies hold working memories in
SK currents: the enrichment marks where to look. The delays here last under a second, shorter than the visual memories of
seconds that depend on signaling in ring neurons of the fly's central complex\cellcite{neuser2008,kuntz2017}, so the two
mechanisms may serve different time scales.

No circuit in the model is necessary for the memory. Its strongest ensemble is densely wired yet no more necessary than matched
sets, and removing all Kenyon cells, the intrinsic cells of the mushroom body, leaves visual memory intact. Flies without
mushroom bodies likewise keep visual learning\cellcite{wolf1998}, although mushroom-body circuits can also store visual
memories\cellcite{vogt2014}. Engrams of a single memory are spread across many brain regions in mice\cellcite{roy2022}; in the
model, the spread follows from where the trace is written and read, and the cue concentrates in the readout cells, several times
out of proportion to their number. Excitability biases which neurons join an engram\cellcite{delamare2024,hansel2024}, and
adaptation currents carried by SK channels may set that bias for a silent working memory.

The model also bears on when to treat. A disease-like shift after learning rewrites the trace and changes its reading, so
normalizing excitability late is a second shift, one that removes the access learned under the disease. Restoring excitability
should therefore help new learning more than the recall of memories formed during the disease. This reading is consistent with
the mixed effects of lowering hyperexcitability in patients\cellcite{bakker2012,vossel2021,faini2025}, with the recovery of
memories by direct engram activation in early disease models\cellcite{roy2016} and with the rescue of retrieval by activating
parvalbumin interneurons in such mice\cellcite{park2024}. Because hyperexcitability emerges long before
diagnosis\cellcite{li2024,okechukwu2026}, a progressing disease keeps moving memories away from the state in which they were
formed, and treatments that hold the state steady early may protect access better than treatments that reverse it late.

Connectome models have so far predicted pathways and activity without learning, or learned within one sensory
pathway\cellcite{shiu2024,lappalainen2024}. A digital twin trained on tasks allows experiments that no animal allows, such as
moving a stored trace from one brain state into another, and each of its results is a prediction for the fly. Because every
experiment here was declared before it ran and replays bit for bit, others can re-examine the twin with the same certainty.
Coupled to a physics model of the fly's body\cellcite{vaxenburg2025}, a trained twin could also link memory to movement in a
closed loop.

Six of these predictions can be tested directly. Knocking down \gene{SK} in carrier hemilineages (0A, 13A/19A, 6A and 21A),
which driver lines can target\cellcite{harris2015}, should shorten visual delay memory more than knocking it down in
\gene{SK}-poor ones (9B, 14A and 10B). An impulse during the delay should reveal the cue in the cell types the model
marks\cellcite{wolff2017}. Training across sustained excitability shifts too small to change behavior should protect later
recall, whereas a shift confined to encoding, or to recall, should impair short visual memory. With temporal control of
transgene expression\cellcite{mcguire2003}, memories formed before Aβ42 is induced should be lost while new learning
works\cellcite{ping2015}, and acute restoration of Kv4 after training under Aβ42 should impair that memory. Memory failure, in
this model, need not mean lost storage: it can mean a key that no longer fits its lock.
